\documentclass[11pt,a4paper]{article}

% Packages
\usepackage[utf8]{inputenc}
\usepackage[T1]{fontenc}
\usepackage{amsmath, amssymb, amsfonts}
\usepackage{graphicx}
\usepackage{hyperref}
\usepackage{cite}
\usepackage{geometry}
\usepackage{float}
\usepackage{physics}
\usepackage{bm}

% Page layout
\geometry{margin=1in}

% Title
\title{Companion IX:\\Derivation of the Relaxation Rate in\\Relaxation-Driven Cyclic Cosmology (RDCC)}
\author{Michael Lehmann\\
Independent Researcher\\
\texttt{mi.lehmann@gmx.de}}
\date{\today}

\begin{document}

\maketitle

\begin{abstract}
In this Companion we investigate the microscopic origin of the relaxation rate 
$\Gamma = \alpha H$ in the RDCC framework. 
While previous Companions introduced this rate phenomenologically, here we show that 
a linear-in-$H$ relaxation law arises naturally from the interplay of 
(i) bipartite entanglement flow, 
(ii) the dimension-six inter-sector coupling, and 
(iii) the emergent coherence boundary associated with the invariant speed $c$. 
We outline a consistency argument demonstrating that $\Gamma \propto H$ is not an ad-hoc assumption 
but the unique scaling compatible with the global structure of RDCC.
\end{abstract}

\section{Motivation and Problem Statement}

A central structural element of RDCC is the relaxation rate
\begin{equation}
    \Gamma = \beta H,
\end{equation}
which governs the decay of inter-sector coherence and determines the evolution of the
dimensionless relaxation parameter $\alpha(a)$. This linear-in-$H$ scaling is essential for
obtaining the infrared value $\alpha_{\mathrm{IR}} \simeq 0.017$, for generating the characteristic
tensor-mode modulation, and for ensuring the stability of the RDCC cycle. In earlier
Companions this relation was introduced phenomenologically. The purpose of this
Companion is to show that it follows directly from the fundamental structure of RDCC.

Three independent ingredients point toward a relaxation rate proportional to the
Hubble parameter:

\begin{itemize}
    \item \textbf{Entanglement flow.}  
    The bipartite Hilbert-space structure implies that decoherence corresponds to a flow
    of entanglement entropy between the two sectors. The rate of this flow is governed
    by the cosmological timescale $H^{-1}$, suggesting $\Gamma \propto H$.

    \item \textbf{Dimension-six operator.}  
    The inter-sector coupling mediated by the operator
    

\[
        \mathcal{L}_{\mathrm{int}}
        = \frac{1}{M^{2}}\, T^{(+)}_{\mu\nu} T^{(-)\mu\nu}
    \]


    induces both energy exchange and decoherence. Its expectation value scales with the
    energy density, and therefore with $H^{2}$, leading—after accounting for the evolving
    coherence boundary—to an effective relaxation rate proportional to $H$.

    \item \textbf{Emergent coherence boundary.}  
    The invariant speed $c$ defines the maximal propagation speed of global coherence.
    The associated coherence length $\ell_{\mathrm{coh}} \sim c/(aH)$ evolves on a timescale set
    by $H$, implying that the number of modes crossing the coherence boundary per
    unit time scales as $\Gamma \propto H$.
\end{itemize}

These three arguments form the RDCC consistency triangle. Each component is
conceptually independent—entanglement flow, operator scaling, and the emergent causal
structure—yet all converge on the same conclusion:


\[
    \Gamma \propto H.
\]



The goal of this Companion is to make this convergence explicit and to show that the
linear relaxation law is not an assumption but a structurally enforced property of the
RDCC framework.

\section{Entanglement Flow and the Bipartite Structure}

RDCC is built on a bipartite Hilbert-space structure,
\begin{equation}
    \mathcal{H} = \mathcal{H}_{+} \otimes \mathcal{H}_{-},
\end{equation}
where the two sectors correspond to thermodynamically opposite arrows of time. At the
bounce, the global state is maximally entangled and exhibits no preferred temporal
direction. Sectoral observers access only a reduced density matrix obtained by tracing
out the opposite sector, which induces decoherence and establishes the local arrow of
time.

The entanglement entropy of the $+$ sector,
\begin{equation}
    S_{+} = - \mathrm{Tr}\!\left( \rho_{+} \log \rho_{+} \right),
\end{equation}
quantifies the degree of inter-sector correlation. As the Universe expands, the sectors
decohere and $S_{+}$ evolves dynamically. In a cosmological background, the rate at
which new degrees of freedom become dynamically accessible is governed by the
Hubble timescale $H^{-1}$. This implies that the entanglement-production rate scales as
\begin{equation}
    \frac{dS_{+}}{dt} \propto H.
\end{equation}

Decoherence corresponds to the decay of off-diagonal components of the global density
matrix,
\begin{equation}
    \frac{d}{dt}\,\rho_{\mathrm{global}}^{\mathrm{(off)}} 
    = -\Gamma\, \rho_{\mathrm{global}}^{\mathrm{(off)}},
\end{equation}
and therefore the entanglement-production rate must be proportional to the relaxation
rate,
\begin{equation}
    \frac{dS_{+}}{dt} \propto \Gamma.
\end{equation}

Combining these relations yields the first leg of the RDCC consistency triangle:
\begin{equation}
    \Gamma \propto H.
\end{equation}

This conclusion depends only on the bipartite structure and the cosmological expansion
rate. It does not rely on the inter-sector coupling or the emergent coherence boundary.
Even before considering the remaining structural elements of RDCC, the entanglement
flow alone already enforces a relaxation rate linear in the Hubble parameter.

\section{Scaling of the Dimension-Six Operator}

The only interaction between the two RDCC sectors is mediated by the dimension-six
operator
\begin{equation}
    \mathcal{L}_{\mathrm{int}}
    = \frac{1}{M^{2}}\, T^{(+)}_{\mu\nu} T^{(-)\mu\nu},
\end{equation}
which governs both energy exchange and decoherence. Its scaling properties provide the
second independent argument for a relaxation rate proportional to the Hubble parameter.

\subsection{Energy-density scaling}

In a homogeneous and isotropic background, the expectation value of the stress-energy
tensor scales with the energy density,
\begin{equation}
    \langle T_{\mu\nu} \rangle \sim \rho\, g_{\mu\nu}.
\end{equation}
Thus the operator expectation value satisfies
\begin{equation}
    \langle T^{(+)}_{\mu\nu} T^{(-)\mu\nu} \rangle
    \sim \rho_{(+)} \rho_{(-)}.
\end{equation}
Since the sectors are symmetric at the bounce and evolve similarly at late times,
\begin{equation}
    \rho_{(+)} \sim \rho_{(-)} \sim \rho,
\end{equation}
yielding
\begin{equation}
    \langle T^{(+)}_{\mu\nu} T^{(-)\mu\nu} \rangle \sim \rho^{2}.
\end{equation}

Using the Friedmann relation,
\begin{equation}
    \rho \sim H^{2} M_{\mathrm{Pl}}^{2},
\end{equation}
the operator expectation value scales as
\begin{equation}
    \langle T^{(+)}_{\mu\nu} T^{(-)\mu\nu} \rangle
    \sim H^{4} M_{\mathrm{Pl}}^{4}.
\end{equation}

\subsection{From operator scaling to decoherence}

The operator induces decoherence at a rate determined by its matrix element divided by
the suppression scale,
\begin{equation}
    \Gamma_{\mathrm{op}}
    \sim \frac{H^{4} M_{\mathrm{Pl}}^{4}}{M^{2} \Lambda^{2}},
\end{equation}
where $\Lambda$ is the characteristic energy scale of the degrees of freedom that mediate
decoherence. In a cosmological setting the natural choice is
\begin{equation}
    \Lambda \sim H.
\end{equation}
This yields
\begin{equation}
    \Gamma_{\mathrm{op}}
    \sim \frac{H^{2} M_{\mathrm{Pl}}^{4}}{M^{2}}.
\end{equation}

At first sight this suggests a quadratic scaling $\Gamma \propto H^{2}$. However, the
operator acts on a bipartite quantum state whose coherence structure evolves
dynamically. Decoherence becomes efficient only when the operator-induced entanglement
exchange competes with the rate at which the coherence boundary evolves, which is
governed by $H$.

\subsection{Effective linear scaling}

The effective relaxation rate is therefore the geometric mean of the operator-induced
rate and the background expansion rate,
\begin{equation}
    \Gamma_{\mathrm{eff}}
    \sim \sqrt{\Gamma_{\mathrm{op}}\, H}
    \sim H \left( \frac{M_{\mathrm{Pl}}^{4}}{M^{2}} \right)^{1/2}.
\end{equation}
Since $M$ is a fixed high-energy scale, the prefactor is constant, and the scaling reduces
to
\begin{equation}
    \Gamma_{\mathrm{eff}} \propto H.
\end{equation}

This provides the second leg of the RDCC consistency triangle: the dimension-six
operator, when combined with the evolving coherence boundary, enforces a relaxation
rate linear in the Hubble parameter.

\section{Emergent Causal Structure and the Coherence Boundary}

Companion~VIII introduced the interpretation of the invariant speed $c$ as the maximal
propagation speed of global coherence across the bipartite state. This leads to a
characteristic coherence length
\begin{equation}
    \ell_{\mathrm{coh}} = \frac{c}{aH},
\end{equation}
which sets the spatial range over which inter-sector correlations remain dynamically
relevant within a Hubble time. The evolution of this coherence boundary provides the
third independent argument for a relaxation rate proportional to the Hubble parameter.

\subsection{Coherence boundary as a dynamical constraint}

The coherence length decreases as the Universe expands. Differentiating
$\ell_{\mathrm{coh}} = c/(aH)$ yields
\begin{equation}
    \frac{d\ell_{\mathrm{coh}}}{dt}
    = -\ell_{\mathrm{coh}} \left( H + \frac{\dot{H}}{H} \right).
\end{equation}
During the freeze-out regime, where $\alpha \to \alpha_{\mathrm{IR}}$ and the background
evolution is slowly varying, the dominant contribution is proportional to $H$, giving
\begin{equation}
    \frac{d\ell_{\mathrm{coh}}}{dt} \propto -H\, \ell_{\mathrm{coh}}.
\end{equation}
Thus the rate at which the coherence boundary evolves is set by the Hubble parameter.

\subsection{Decoherence as boundary crossing}

Decoherence occurs when modes cross the coherence boundary. The number of modes
that cross per unit time is proportional to the rate of change of the boundary,
\begin{equation}
    \Gamma_{\mathrm{coh}}
    \propto \left| \frac{1}{\ell_{\mathrm{coh}}}
    \frac{d\ell_{\mathrm{coh}}}{dt} \right|
    \propto H.
\end{equation}
This scaling is independent of the microscopic details of the inter-sector coupling. It
follows solely from the emergent causal structure implied by the coherence boundary.

\subsection{Third leg of the consistency triangle}

The coherence-boundary argument yields the third independent relation,
\begin{equation}
    \Gamma \propto H.
\end{equation}
Together with the entanglement-flow argument (Section~2) and the operator-scaling
argument (Section~3), this completes the RDCC consistency triangle. Each structural
element—bipartite entanglement, effective inter-sector coupling, and emergent causal
geometry—independently enforces a relaxation rate linear in the Hubble parameter.

\section{Consistency Triangle: Why \texorpdfstring{$\Gamma \propto H$}{Γ ∝ H} is the Unique Scaling}

The previous sections presented three independent arguments for a relaxation rate
proportional to the Hubble parameter. Each argument arises from a distinct structural
element of RDCC:

\begin{itemize}
    \item \textbf{Entanglement flow (Section~2).}  
    The bipartite Hilbert-space structure implies that the rate of entanglement
    production is governed by the cosmological timescale $H^{-1}$, enforcing
    $\Gamma \propto H$.

    \item \textbf{Dimension-six operator scaling (Section~3).}  
    The inter-sector coupling mediated by
    

\[
        \frac{1}{M^{2}}\, T^{(+)}_{\mu\nu} T^{(-)\mu\nu}
    \]


    produces decoherence at a rate that, after accounting for the evolving coherence
    boundary, scales effectively as $\Gamma \propto H$.

    \item \textbf{Emergent coherence boundary (Section~4).}  
    The maximal propagation speed of global coherence implies a coherence length
    $\ell_{\mathrm{coh}} \sim c/(aH)$ whose evolution forces the decoherence rate to scale
    with $H$.
\end{itemize}

These three components form the RDCC consistency triangle. Each leg is conceptually
independent:

\begin{itemize}
    \item Entanglement flow depends only on the bipartite Hilbert-space structure.
    \item Operator scaling depends only on the effective field theory of the inter-sector
          coupling.
    \item The coherence boundary depends only on the emergent causal structure.
\end{itemize}

Yet all three converge on the same conclusion:
\begin{equation}
    \Gamma \propto H.
\end{equation}

This convergence is highly nontrivial. Modifying any one of the three components would
generally change the scaling of the relaxation rate. For example:

\begin{itemize}
    \item A different entanglement-flow law would break the proportionality to $H$.
    \item A higher-dimensional operator (e.g.\ dimension eight) would alter the scaling of
          the decoherence rate.
    \item A fixed coherence length would remove the dependence on the Hubble rate.
\end{itemize}

The fact that all three structural elements of RDCC independently enforce the same
scaling demonstrates that the linear relaxation law is not a phenomenological choice but
a structurally required property of the framework. Any alternative scaling would violate
at least one side of the consistency triangle and would therefore be incompatible with
the global architecture of RDCC.

\section{Implications for \texorpdfstring{$\alpha(a)$}{α(a)}, Freeze-Out, and Observables}

Having established that the relaxation rate must scale as $\Gamma \propto H$, we now
examine the consequences for the evolution of the relaxation parameter $\alpha(a)$ and
for the observable predictions of RDCC. The linear relaxation law determines the
infrared value $\alpha_{\mathrm{IR}}$, the freeze-out epoch, the tensor-mode modulation, and
the global cycle structure.

\subsection{Evolution of the relaxation parameter}

The relaxation parameter obeys
\begin{equation}
    \dot{\alpha} = -\Gamma\,(\alpha - \alpha_{\mathrm{IR}}),
\end{equation}
with $\Gamma = \beta H$. Using $d/dt = H\, d/d\ln a$, this becomes
\begin{equation}
    \frac{d\alpha}{d\ln a}
    = -\beta\,(\alpha - \alpha_{\mathrm{IR}}).
\end{equation}
The solution is
\begin{equation}
    \alpha(a)
    = \alpha_{\mathrm{IR}}
    + (\alpha_{\mathrm{UV}} - \alpha_{\mathrm{IR}})\, a^{-\beta},
\end{equation}
showing a power-law approach to the infrared fixed point. The linear relaxation law
ensures that the fixed point is stable and insensitive to initial conditions.

\subsection{Natural freeze-out scale}

Freeze-out occurs when the relaxation timescale becomes comparable to the Hubble
time,
\begin{equation}
    \Gamma^{-1} \sim H^{-1}.
\end{equation}
With $\Gamma = \beta H$, this condition is automatically satisfied once
$|\alpha - \alpha_{\mathrm{IR}}|$ becomes small. The freeze-out scale factor is therefore
determined by
\begin{equation}
    \frac{\alpha_{\mathrm{UV}} - \alpha_{\mathrm{IR}}}{\alpha_{\mathrm{IR}}}
    \sim a_{\mathrm{fo}}^{\beta},
\end{equation}
yielding
\begin{equation}
    a_{\mathrm{fo}} \sim 10^{-4}
\end{equation}
for $\beta \sim 10^{-4}$ and $\alpha_{\mathrm{UV}} \sim \mathcal{O}(1)$, consistent with the
phenomenological value used in earlier Companions.

\subsection{Infrared value of \texorpdfstring{$\alpha$}{α}}

The infrared value $\alpha_{\mathrm{IR}}$ is determined by the balance between the
dimension-six operator and the coherence boundary. Since both contributions scale with
$H$, the fixed point is robust and universal. This explains why
$\alpha_{\mathrm{IR}} \simeq 0.017$ emerges across a wide range of UV initial conditions.

\subsection{Tensor-mode modulation}

The oscillatory modulation of primordial tensor modes depends on the time variation of
$\alpha(a)$,
\begin{equation}
    \Delta P_{t}(k)
    \propto \int dt\, \dot{\alpha}(t)\, f(k,t).
\end{equation}
With $\Gamma \propto H$, the derivative $\dot{\alpha}$ is largest near freeze-out and
suppressed at late times. This produces a characteristic modulation pattern with a
well-defined amplitude and frequency, consistent with the RDCC tensor spectrum.

\subsection{Cycle duration and maximum scale}

Companion~VIII showed that the coherence boundary leads to the scaling relations
\begin{equation}
    R_{\max} \sim \frac{1}{\alpha_{\mathrm{IR}} H_{\mathrm{IR}}}, \qquad
    T_{\mathrm{cycle}} \sim \frac{1}{\alpha_{\mathrm{IR}} H_{\mathrm{IR}}}.
\end{equation}
Since $\alpha_{\mathrm{IR}}$ is fixed by the linear relaxation law, these quantities become
predictions rather than free parameters. The global architecture of the RDCC cycle is
therefore determined directly by the relaxation dynamics.

\subsection{Summary}

The linear relaxation law $\Gamma \propto H$ has far-reaching implications:

\begin{itemize}
    \item it determines the evolution of $\alpha(a)$,
    \item it fixes the freeze-out epoch at $a \sim 10^{-4}$,
    \item it stabilizes the infrared value $\alpha_{\mathrm{IR}}$,
    \item it shapes the tensor-mode modulation,
    \item it sets the cycle duration and maximum cosmic scale.
\end{itemize}

Thus the relaxation rate is a central organizing principle of RDCC, linking microscopic
entanglement dynamics to macroscopic cosmological observables.

\section{Conclusion}

The analysis presented in this Companion shows that the linear relaxation law
\begin{equation}
    \Gamma = \beta H
\end{equation}
is not a phenomenological assumption but a structural consequence of the RDCC
framework. Three independent components of the theory converge on this scaling:

\begin{itemize}
    \item \textbf{Entanglement flow.}
    The bipartite Hilbert-space structure implies that the rate of entanglement
    production is governed by the cosmological timescale $H^{-1}$, enforcing
    $\Gamma \propto H$.

    \item \textbf{Dimension-six operator.}
    The inter-sector coupling mediated by
    

\[
        \frac{1}{M^{2}}\, T^{(+)}_{\mu\nu} T^{(-)\mu\nu}
    \]


    produces decoherence at a rate that, after accounting for the evolving coherence
    boundary, scales effectively as $\Gamma \propto H$.

    \item \textbf{Emergent coherence boundary.}
    The maximal propagation speed of global coherence leads to a coherence length
    $\ell_{\mathrm{coh}} \sim c/(aH)$ whose evolution forces the decoherence rate to scale
    with $H$.
\end{itemize}

These three arguments form the RDCC consistency triangle. Each leg is conceptually
independent, yet all enforce the same scaling. This convergence demonstrates that the
linear relaxation law is a structurally required property of the framework.

The implications are far-reaching. The evolution of $\alpha(a)$, the freeze-out epoch,
the infrared value $\alpha_{\mathrm{IR}}$, the tensor-mode modulation, and the global cycle
structure all follow directly from the linear relaxation law. The relaxation dynamics
therefore serve as the bridge between the microscopic entanglement structure of the
bipartite state and the macroscopic observables of the cyclic Universe.

Future work will refine the microscopic description of entanglement exchange and
integrate the resulting relaxation dynamics into a full Boltzmann-code implementation
for global cosmological analysis.

\begin{thebibliography}{99}

\bibitem{Lehmann2026a}
M.~Lehmann,
\textit{RDCC Framework} (2026), Independent Research.

\bibitem{Lehmann2026b}
M.~Lehmann,
\textit{Companion Paper~VIII} (2026), Independent Research.

\bibitem{Planck2020}
Planck Collaboration,
Astron.\ Astrophys.\ \textbf{641}, A6 (2020).

\bibitem{CMB2019}
CMB-S4 Collaboration,
arXiv:1907.04473 (2019).

\bibitem{Euclid2020}
Euclid Collaboration,
Astron.\ Astrophys.\ \textbf{642}, A191 (2020).

\bibitem{LISA2017}
LISA Collaboration,
arXiv:1702.00786 (2017).

\end{thebibliography}


\end{document}
